\section*{第 \ref{ch:b2:quadratic} 章 --- 二次型}

\begin{solution}{exo:b2:quadratic:1}
$q_1 = (x + 2y)^2 - 3y^2$：秩为 $2$，符号差为 $(1, 1)$（
双曲型二次型）。

$q_2$：对 $x$ 配方：$q_2 = (x + y)^2 + y^2 + 2yz +
3z^2 = (x+y)^2 + (y + z)^2 + 2z^2$：秩为 $3$，符号差为 $(3, 0)$
--- 正定。
\end{solution}

\begin{solution}{exo:b2:quadratic:2}
特征值为 $2$（在 $\operatorname{Vect}(1,1,1)$ 上）与 $-1$（在
平面 $x + y + z = 0$ 上）。标准正交化：$e_1 =
\frac{1}{\sqrt3}(1,1,1)$；在该平面中对 $(1,-1,0),
(1,0,-1)$ 施行格拉姆–施密特正交化，得 $e_2 = \frac{1}{\sqrt2}(1,-1,0)$，$e_3 =
\frac{1}{\sqrt6}(1,1,-2)$。于是 $P = (e_1\ e_2\ e_3)$ 为正交矩阵，
且 $P^{\mathsf T}AP = \operatorname{diag}(2,-1,-1)$。

二次型 $q = 2xy + 2yz + 2zx$ 的矩阵为 $A$。在旋转后的坐标系中
$q = 2X^2 - Y^2 - Z^2$ --- 即主轴；符号差
为 $(1,2)$，与\cref{ex:b2:quadratic:gaussexample}
相符（同一二次型！）。
\end{solution}

\begin{solution}{exo:b2:quadratic:3}
$u^{**} = u$：$\langle u^{**}x, y\rangle = \langle x, u^*y\rangle =
\langle ux, y\rangle$ 对所有 $y$ 成立。$(uv)^* = v^*u^*$：$\langle
uvx, y\rangle = \langle vx, u^*y\rangle = \langle x,
v^*u^*y\rangle$。核：$y \in \ker u^* \iff \langle x, u^*y
\rangle = 0\ \forall x \iff \langle u(x), y\rangle = 0\ \forall x
\iff y \perp \operatorname{im} u$。秩：$\dim\ker u^* = n -
\operatorname{rk} u$（正交补），故
$\operatorname{rk} u^* = \operatorname{rk} u$ 由秩–零化度定理得 ---
这是欧氏空间中转置秩定理的版本。
\end{solution}

\begin{solution}{exo:b2:quadratic:4}
谱分解：$A = PDP^{\mathsf T}$，$D$ 为对角矩阵，其元素
$\lambda_i$ 满足 $\lambda_i^3 = \lambda_i$，故 $\lambda_i \in
\{-1, 0, 1\}$。于是 $A^2 = PD^2P^{\mathsf T}$，其中 $D^2$ 为对角矩阵，
其元素为 $0/1$：$A^2$ 对称且幂等
（$(A^2)^2 = A^4 = A\cdot A^3 = A^2$） --- 对称幂等矩阵
$=$ 正交投影（即到
$\ker(A^2 - I) = \ker(A-I)\oplus\ker(A+I)$ 上、沿 $\ker A$ 的投影，
二者由谱定理知正交）。

一般地，$P(A) = P\!\left(\text{diag}\right)$：$P(A)$ 具有
相同的特征向量，特征值为 $P(\lambda_i)$ --- 这就是
可对角化情形下的“谱映射”。
\end{solution}

\begin{solution}{exo:b2:quadratic:5}
闭：$O(n) = g^{-1}(\{I\})$，其中连续映射 $g(P) =
P^{\mathsf T}P$（元素为多项式）。有界：$P
\in O(n)$ 的每一列都是单位向量，故其所有元素位于 $\intcc{-1}{1}$ 中。
在 $\mathcal{M}_n(\R) \simeq \R^{n^2}$ 中既闭又有界，故
为紧集（\cref{thm:b2:metric:compactprops}~(2)）。

不连通：$\det$ 取两个值 $\pm1$，二者均在 $O(n)$ 上取得，
而连续满射到 $\{-1, 1\}$ 上将空间分裂
（即\cref{ex:b2:metric:glnr}中的论证）。
\end{solution}

\begin{solution}{exo:b2:quadratic:6}
\emph{存在性：}$A = PDP^{\mathsf T}$，其中 $D =
\operatorname{diag}(\lambda_i)$，$\lambda_i \geq 0$；令 $B =
P\sqrt D P^{\mathsf T}$，其中 $\sqrt D =
\operatorname{diag}(\sqrt{\lambda_i})$：对称且
半正定，$B^2 = A$。

\emph{唯一性：}设 $B$ 为对称半正定矩阵，且 $B^2 = A$。$B$
与 $A$ 可交换；因此 $B$ 保持每个特征子空间
$E_\lambda(A)$ 不变（对 $Ax = \lambda x$，有 $A(Bx) = BAx = \lambda Bx$）。
在 $E_\lambda(A)$ 上，$B$ 的限制为对称半正定，
其平方为 $\lambda\,\mathrm{id}$；其特征值 $\mu$ 满足 $\mu^2
= \lambda$，$\mu \geq 0$，于是 $\mu = \sqrt\lambda$ --- 因此
该限制可对角化，且唯一特征值为
$\sqrt\lambda$，故它\emph{就是} $\sqrt\lambda\,\mathrm{id}$。由 $E =
\bigoplus E_\lambda(A)$，$B$ 被唯一确定：$B = \sqrt A$。
\end{solution}

\begin{solution}{exo:b2:quadratic:7}
在标准正交特征基下，$\norm{Ax}_2^2 = \sum \lambda_i^2
x_i^2 \leq (\max_i \lambda_i^2)\norm x_2^2$，等号在相应的
特征向量处取得：$\vertiii A_2 = \max\abs{\lambda_i}$。
对所给矩阵，特征值为 $3, -1$
（与\cref{ex:b2:quadratic:spectralexample}相似）：
$\vertiii A_2 = 3$。
\end{solution}

\begin{solution}{exo:b2:quadratic:8}
（$\Rightarrow$）左上角 $k \times k$ 子块 $A_k$ 是正定二次型
限制到前 $k$ 个基向量张成的子空间上的矩阵，
即为正定矩阵，故其特征值均为正，
从而 $\Delta_k = \det A_k > 0$。

（$\Leftarrow$）对 $n$ 用归纳法；$n = 1$ 时显然。设全部
$\Delta_k > 0$。由归纳假设，$A_{n-1}$ 正定，即
二次型 $q$ 限制到 $F = \operatorname{Vect}(e_1, \dots,
e_{n-1})$ 上正定。用高斯法将 $q|_F$ 对角化，取
坐标 $y_1, \dots, y_{n-1}$，使 $q|_F = \sum y_i^2$。在整个空间中，
对最后一个变量配方，
\[
q = \sum_{i=1}^{n-1} \bigl(y_i + c_i x_n\bigr)^2 + c\,x_n^2
\]
其中常数适当选取（将交叉项并入各平方项）。
该化简给出符号差 $(n-1 + \epsilon, \cdot)$，其中
$\epsilon$ 为 $c$ 的符号贡献；而行列式
在合同下保持对角系数之积的符号
（$\det(P^{\mathsf T}AP) = (\det P)^2\det A$）：$\Delta_n
> 0$ 迫使 $c > 0$。因此 $q$ 是 $n$ 个
线性无关线性型的平方和：正定。
\end{solution}

\begin{solution}{exo:b2:quadratic:9}
设 $(e_1, \dots, e_n)$ 为对应于 $\lambda_1
\geq \dots \geq \lambda_n$ 的一组标准正交特征基。

\emph{$\lambda_2 \leq$ 极小极大值：}对任意超平面 $H$，
$2$ 维子空间 $V = \operatorname{Vect}(e_1, e_2)$ 满足
由格拉斯曼公式，$\dim(H \cap V) \geq 1$，取单位向量 $x \in H \cap V$，
写 $x = ae_1 + be_2$，$a^2 + b^2 = 1$：
\[
\langle u(x), x\rangle = \lambda_1 a^2 + \lambda_2 b^2 \geq
\lambda_2 :
\]
于是每个超平面的最大值 $\geq \lambda_2$。

\emph{$\geq$：}对 $H = e_1^{\perp}$，任一单位向量 $x = \sum_{i\geq2}
x_ie_i \in H$ 满足 $\langle u(x), x\rangle = \sum_{i \geq 2}
\lambda_i x_i^2 \leq \lambda_2$，在 $e_2$ 处取到：此
超平面上的最大值恰为 $\lambda_2$。对 $H$ 取极小，
因此其最小值为 $\lambda_2$。
\end{solution}

\begin{solution}{exo:b2:quadratic:10}
\emph{代数途径：}$2q = \bigl(\sum x_i\bigr)^2 - \sum
x_i^2$。在超平面 $H : \sum x_i = 0$ 上（维数为 $n -
1$），$q = -\frac12\sum x_i^2$ 为负定；在直线
$\R(1, \dots, 1)$ 上，$q(t, \dots, t) = \binom n2 t^2 > 0$。
使 $q$ 正定的子空间与 $H$ 仅交于零，故
其维数 $\leq 1$：由西尔维斯特
（\cref{thm:b2:quadratic:sylvester}）知 $s = 1$，且 $t \geq n-1$（由 $H$ 得）；
秩 $\leq n$ 迫使符号差为 $(1, n-1)$，秩为 $n$。

\emph{谱途径：}矩阵为 $\frac12(J - I)$；$J$ 有
特征值 $n$（对应 $(1,\dots,1)$）与 $0$（对应 $H$），故
$\frac12(J-I)$ 的特征值为 $\frac{n-1}{2}$（单重）与
$-\frac12$（$n-1$ 重）：正的一个，负的 $n-1$ 个 --- 符号差
相同，由\cref{cor:b2:quadratic:principalaxes}可知。
\end{solution}

\begin{solution}{exo:b2:quadratic:11}
记 $B = C^{\mathsf T}C$（见\cref{exo:b2:quadratic:6}，取 $C =
\sqrt B$）。于是记 $c_1, \dots, c_n$ 为
$C^{\mathsf T}$ 的各列，则：
\[
\operatorname{tr}(AB) = \operatorname{tr}(AC^{\mathsf T}C)
= \operatorname{tr}(CAC^{\mathsf T})
= \sum_{i=1}^n \langle A c_i, c_i\rangle .
\]
由\cref{cor:b2:quadratic:principalaxes}~(2)，每一项
介于 $\lambda_{\min}(A)\norm{c_i}^2$ 与
$\lambda_{\max}(A)\norm{c_i}^2$ 之间，且 $\sum\norm{c_i}^2 =
\operatorname{tr}(C^{\mathsf T}C)^{\vphantom1} =
\operatorname{tr} B$，于是双边不等式成立。若 $A$
也是半正定的，则 $\lambda_{\min}(A) \geq 0$：
$\operatorname{tr}(AB) \geq 0$。
\end{solution}

\begin{solution}{exo:b2:quadratic:12}
\begin{enumerate}
  \item $\varphi(M, N) = \operatorname{tr}(MN)$ 是双线性的，且
        对称（$\operatorname{tr}(MN) =
        \operatorname{tr}(NM)$），且 $\varphi(M, M) = q(M)$：
        $q$ 是 $\varphi$ 的二次型。
  \item 设 $S$ 对称、$K$ 反对称：
        $\operatorname{tr}(SK) =
        \operatorname{tr}\bigl((SK)^{\mathsf T}\bigr) =
        \operatorname{tr}(K^{\mathsf T}S^{\mathsf T}) =
        -\operatorname{tr}(KS) = -\operatorname{tr}(SK)$，故
        $\varphi(S, K) = 0$：两子空间是
        $\varphi$ 正交的。按元素展开，$\operatorname{tr}
        (M^2) = \sum_{i,j} m_{ij}m_{ji}$：当 $M$ 对称时
        此式为 $\sum m_{ij}^2 > 0$（$M \neq 0$）；对
        反对称的 $M$，则为 $-\sum m_{ij}^2 < 0$。
  \item $\mathcal M_n(\R) = S_n \oplus A_n$，维数分别为
        $\frac{n(n+1)}2$ 与 $\frac{n(n-1)}2$；取一个高斯化简，
        使其适应此 $\varphi$ 正交分解，
        由此分解将 $q$ 写成 $\frac{n(n+1)}2$ 个正平方项与
        $\frac{n(n-1)}2$ 个负平方项：符号差
        $\bigl(\frac{n(n+1)}2, \frac{n(n-1)}2\bigr)$
        （西尔维斯特），秩为 $n^2$：该型非退化。
\end{enumerate}
\end{solution}

\begin{solution}{pb:b2:quadratic:1}
\textbf{1.} 由内积的对称性，$G$ 是对称的，
并有
\[
X^{\mathsf T}GX = \sum_{i,j}x_ix_j\langle v_i, v_j\rangle
= \Bigl\|\sum_i x_iv_i\Bigr\|^2 \geq 0 ,
\]
等号成立当且仅当 $\sum x_iv_i = 0$：$G$ 正定，当且仅当
仅有平凡的零组合，即该向量族
线性无关。

\textbf{2.} 取 $B = \sqrt A$（见\cref{exo:b2:quadratic:6}）：$A
= B^2 = B^{\mathsf T}B$，即 $B$ 各列的格拉姆矩阵；
取 $C = B$。又 $X^{\mathsf T}AX = \norm{CX}^2$，故 $A$
正定，当且仅当 $CX \neq 0$ 对一切 $X \neq 0$ 成立，即 $C$ 可逆。

\textbf{3.} 将该型限制到
$\operatorname{Vect}(e_i, e_j)$ 上，其矩阵为
$\begin{pmatrix} a_{ii} & a_{ij}\\ a_{ij} & a_{jj}
\end{pmatrix}$，仍为半正定，其行列式
（为其非负特征值之积）$\geq 0$：
$a_{ij}^2 \leq a_{ii}a_{jj}$。此即
对格拉姆实现中向量 $v_i, v_j$ 的柯西–施瓦茨不等式。

\textbf{4.} \emph{存在性：}将 $A$ 写为
线性无关向量组 $(v_1, \dots, v_n)$ 的格拉姆矩阵（由第 1–2 问）。
格拉姆–施密特过程给出标准正交组 $(e_1, \dots, e_n)$，使
\[
v_k = \sum_{i \leq k} t_{ik}\,e_i,
\qquad t_{kk} = \bigl\| v_k - \operatorname{proj}_{k-1}v_k
\bigr\| > 0 ,
\]
于是 $T = (t_{ik})$ 为对角元为正的上三角矩阵，
并有
\[
a_{jk} = \langle v_j, v_k\rangle
= \sum_i t_{ij}t_{ik} = (T^{\mathsf T}T)_{jk} .
\]
\emph{唯一性：}若 $T_1^{\mathsf T}T_1 = T_2^{\mathsf T}T_2$
则 $R = T_1T_2^{-1}$ 满足 $R^{\mathsf T}R = I$：$R$ 为
正交矩阵，且也是对角元为正的上三角矩阵
（此类矩阵之积）。于是 $R^{-1} = R^{\mathsf T}$ 同时
是上三角（上三角矩阵之逆）与下三角（上三角矩阵的
转置）矩阵，故为对角矩阵；正交的对角矩阵
元素为 $\pm1$，而正性迫使 $R = I$：$T_1 = T_2$。

\textbf{5.} 对 $i, j \leq k$，$(T^{\mathsf T}T)_{ij} = \sum_m
t_{mi}t_{mj}$ 只涉及 $m \leq \min(i,j) \leq k$ 的项，因此
$k\times k$ 前导块（取自 $A$）为 $T_k^{\mathsf T}T_k$，其中
$T_k$ 为 $T$ 的前导块。于是 $\Delta_k = (\det T_k)^2
= (t_{11}\cdots t_{kk})^2$。现在，按自然顺序对正定型作高斯约化
从不会遇到零平方系数
（主元即逐次约化所得正定块的对角元），
得到 $q
= \sum_k d_k\ell_k^2$，其中 $\ell_k = x_k + (\text{ 含 }
x_{k+1}, \dots)$ 的项，即 $A = L^{\mathsf T}DL$，
$L$ 为单位上三角矩阵；于是 $T = \sqrt D\,L$ 为一个楚列斯基
分解，由唯一性得 $t_{kk}^2 = d_k$，从而
\[
d_k = \frac{(t_{11}\cdots t_{kk})^2}
{(t_{11}\cdots t_{k-1,k-1})^2}
= \frac{\Delta_k}{\Delta_{k-1}} .
\]
在\cref{ex:b2:quadratic:tworoads}中：$\Delta_1, \Delta_2,
\Delta_3 = 2, 3, 4$ 且主元依次为 $2, \frac32, \frac43$。

\textbf{6.} 每个 $a_{ii} = e_i^{\mathsf T}Ae_i > 0$。令 $D =
\operatorname{diag}(a_{ii}^{-1/2})$，$B = DAD$ 正定（合同），
且 $b_{ii} = 1$，故 $\operatorname{tr}
B = n$。其特征值 $\mu_i > 0$ 满足，由算术–几何平均不等式，
\[
\det B = \prod_i\mu_i
\leq \Bigl(\frac{\sum\mu_i}{n}\Bigr)^{\!n} = 1 ,
\]
又 $\det B = (\det D)^2\det A = \dfrac{\det A}{\prod
a_{ii}}$，于是 $\det A \leq \prod a_{ii}$。

\textbf{7.} 算术–几何平均不等式取等当且仅当所有 $\mu_i$ 相等
（均等于 $1$）；特征值仅为 $1$ 的对称矩阵满足
$PIP^{\mathsf T} = I$。故取等当且仅当 $B = I$，亦即 $a_{ij} = 0$
（$i \neq j$），即 $A$ 为对角矩阵。

\textbf{8.} 若 $M$ 奇异，则两边都 $\geq 0 =
\abs{\det M}$。否则 $A = M^{\mathsf T}M$ 为正定矩阵，
且 $a_{ii} = \norm{c_i}^2$，$\det A = (\det
M)^2$：由第 6 题得 $(\det M)^2 \leq \prod\norm{c_i}^2$。
取等当且仅当 $A = M^{\mathsf T}M$ 为对角矩阵（由第 7 题），
即当且仅当各列两两正交。

\textbf{9.} $\abs{\det M}$ 是由各列张成的平行多面体的体积：
体积不超过各棱长之积，
取等当且仅当为长方体。
若 $\abs{m_{ij}} \leq 1$，则 $\norm{c_i} \leq \sqrt n$，故
$\abs{\det M} \leq n^{n/2}$。（取等时须各列正交，且元素为 $\pm1$，
即阿达马矩阵。）

\textbf{10.} $X^{\mathsf T}A^{\mathsf T}AX = \norm{AX}^2 > 0$，对 $X \neq 0$
（$A$ 可逆），故 $A^{\mathsf T}A$ 正定。其平方根 $S$ 为
正定（特征值
$\sqrt{\lambda_i} > 0$），从而可逆，且 $Q = AS^{-1}$
满足
\[
Q^{\mathsf T}Q = S^{-1}A^{\mathsf T}AS^{-1}
= S^{-1}S^2S^{-1} = I :
\]
$A = QS$，其中 $Q$ 正交，$S$ 正定。

\textbf{11.} 若 $A = QS = Q'S'$，则 $S'^{\,2} =
S'^{\mathsf T}Q'^{\mathsf T}Q'S' = A^{\mathsf T}A = S^2$；两个
平方相同的半正定矩阵必相等
（\cref{exo:b2:quadratic:6}）：$S' = S$，于是 $Q' = AS^{-1} =
Q$.

\textbf{12.} $\det(A + \varepsilon I)$ 是
$\varepsilon$ 的非零多项式，只有有限个根，故存在数列
$\varepsilon_k \to 0$ 避开这些根。记 $A + \varepsilon_kI =
Q_kS_k$（第 10 题）。$O(n)$ 是紧的
（\cref{exo:b2:quadratic:5}）：存在子列 $Q_{\varphi(k)}
\to Q \in O(n)$。于是
\[
S_{\varphi(k)} = Q_{\varphi(k)}^{\mathsf T}
\bigl(A + \varepsilon_{\varphi(k)}I\bigr)
\longrightarrow Q^{\mathsf T}A =: S,
\]
作为此类矩阵的极限，对称半正定（这些条件
为闭的），且 $A = QS$。又 $S^2 = S^{\mathsf T}S =
A^{\mathsf T}QQ^{\mathsf T}A = A^{\mathsf T}A$，故 $S =
\sqrt{A^{\mathsf T}A}$（由唯一性）。若 $A$ 奇异，则 $S$ 奇异，
$Q$ 不唯一：它可在 $(\operatorname{im} S)^{\perp}$ 上任意改变 ---
极端情形 $A = 0$，
此时任一正交矩阵 $Q$ 均可。

\textbf{13.} 将 $S = P\Sigma P^{\mathsf T}$ 对角化（谱定理），
$\Sigma = \operatorname{diag}(\sigma_i)$，
其中 $\sigma_i \geq 0$ 为 $S =
\sqrt{A^{\mathsf T}A}$ 的特征值。于是
\[
A = QS = (QP)\,\Sigma\,P^{\mathsf T} = U\Sigma V^{\mathsf T},
\qquad U = QP,\ V = P \in O(n) .
\]

\textbf{14.} 有 $\norm{Ax}^2 = x^{\mathsf T}S^2x \leq
\sigma_{\max}^2\norm x^2$，等号在 $S$ 的最大特征值
所对应的特征向量处取得：$\vertiii A_2 = \sigma_{\max}$ ---
对于对称矩阵 $A$，$S = \sqrt{A^2}$ 的特征值为 $\abs{\lambda_i}$，
与 \cref{exo:b2:quadratic:7} 相吻合。行列式：$\abs{\det A} =
\abs{\det U}\det\Sigma\abs{\det V} = \sigma_1\cdots\sigma_n$.
球面：令 $x = Vy$，$\norm y = 1$，则 $Ax = U\Sigma y$
的坐标为 $z_i = \sigma_iy_i$，在 $U$ 的列所定的
标准正交标架下，像为 $\{\sum z_i^2/\sigma_i^2 =
1\}$，即半轴长为 $\sigma_i$ 的椭球面。

\textbf{15.} 旋转至主轴
（\cref{cor:b2:quadratic:principalaxes}）后，$f$ 化为
$\lambda_1X^2 + \lambda_2Y^2 + \beta_1X + \beta_2Y + c$，其中
$\lambda_1\lambda_2 = \det A$。若 $\det A \neq 0$，
通过平移 $X \mapsto X -
\frac{\beta_1}{2\lambda_1}$（$Y$ 同理）消去一次项：$\lambda_1X'^2
+ \lambda_2Y'^2 = c'$。若 $\det A > 0$（同号），则为
椭圆（$c'$ 符号适当）、一个点或空集。若 $\det
A < 0$，则为双曲线（$c' \neq 0$）或两条相交直线。若
$\det A = 0$ 且秩为 $1$（设 $\lambda_2 = 0 \neq
\lambda_1$）：$\lambda_1X'^2 + \beta_2Y + c''$。若 $\beta_2 \neq 0$，则为抛物线；
否则为两条平行直线、一条直线或空集。
非退化的图形为椭圆、双曲线、抛物线，
由 $\det A$ 的符号决定。

\textbf{16.} 二次部分 $x^2 + 4xy + y^2$ 的矩阵为
$\begin{pmatrix}1 & 2\\ 2 & 1\end{pmatrix}$，特征值为 $3$
与 $-1$，单位正交的方向分别为 $\frac{1}{\sqrt2}(1,1)$，
$\frac{1}{\sqrt2}(1,-1)$。在旋转后的坐标中，$u =
\frac{x+y}{\sqrt2}$, $v = \frac{x-y}{\sqrt2}$: $x^2 + y^2 = u^2
+ v^2$，$2xy = u^2 - v^2$，于是二次型为 $3u^2 - v^2$，且 $2x
- 2y = 2\sqrt2\,v$。方程化为
\[
3u^2 - v^2 + 2\sqrt2\,v = 4
\quad\Longleftrightarrow\quad
3u^2 - \bigl(v - \sqrt2\bigr)^2 = 2 :
\]
这是双曲线，中心位于 $(u, v) = (0, \sqrt2)$，即 $(x, y) =
(1, -1)$，其轴沿旋转后的坐标系方向。

\textbf{17.} $f(x) = (x - x_0)^{\mathsf T}A(x - x_0) + f(x_0)$
上式在 $Ax_0 = -\frac b2$ 时成立，即 $x_0 = -\frac12A^{-1}b$：
$f$ 的梯度恰在该点为零（$x_0$ 是
对称中心）。取 $M = \begin{pmatrix} I & x_0\\ 0 &
1\end{pmatrix}$:
\[
M^{\mathsf T}\widetilde QM
= \begin{pmatrix}
A & Ax_0 + \frac b2\\[2pt]
\bigl(Ax_0 + \frac b2\bigr)^{\mathsf T} &
x_0^{\mathsf T}Ax_0 + b^{\mathsf T}x_0 + c
\end{pmatrix}
= \begin{pmatrix} A & 0\\ 0 & f(x_0)\end{pmatrix},
\]
且 $\det M = 1$，故 $\det\widetilde Q = f(x_0)\det A$。
中心化方程为 $q(X) = -f(x_0)$：当 $f(x_0) = 0$ 时，它
退化为 $q(X) = 0$（
符号差为 $(1,1)$ 时为过中心的两条直线，单点 $x_0$ 则在 $q$ 定号时出现）；
当 $f(x_0) \neq 0$ 时，二次曲线为真正的椭圆
或双曲线。

\textbf{18.} $A^{-1} = -\frac13\begin{pmatrix} 1 & -2\\ -2 &
1\end{pmatrix}$, $\frac b2 = (1, -1)$: $x_0 =
-A^{-1}\frac b2 = (1, -1)$，如第 16 题所得。$f(x_0) =
q(1,-1) + 2 + 2 - 4 = (1 - 4 + 1) + 0 = -2 \neq 0$，且
$\det\widetilde Q = f(x_0)\det A = (-2)(-3) = 6 \neq 0$:
非退化；$\det A = -3 < 0$：为双曲线 --- 事实上
中心化方程 $3u^2 - (v - \sqrt2)^2 = -f(x_0) = 2$
与第 16 题相符。

\textbf{19.} 矩阵为 $(1-\lambda)I + \lambda J$：
特征值为 $1 + 2\lambda$（方向 $(1,1,1)$）与 $1 -
\lambda$（二重，位于 $x + y + z = 0$ 上）。分情形如下：
\begin{itemize}
  \item $-\frac12 < \lambda < 1$：特征值全为正：
        以 $(1,1,1)$ 为轴的旋转椭球面（当
        $\lambda = 0$ 时为球面）；
  \item $\lambda = 1$：$q = (x+y+z)^2$：方程给出
        两个平行平面 $x + y + z = \pm1$；
  \item $\lambda = -\frac12$：特征值为 $0, \frac32,
        \frac32$：以 $(1,1,1)$ 为轴的圆柱面；
  \item $\lambda > 1$：符号差为 $(1, 2)$，为
        双叶双曲面；
  \item $\lambda < -\frac12$：符号差为 $(2, 1)$，为
        单叶双曲面。
\end{itemize}

\textbf{20.} $q$ 的极化双线性型
$\langle\cdot,\cdot\rangle_q$ 为 $E$ 上的内积。固定 $x$，$y \mapsto
\varphi'(x, y)$（$q'$ 的极化）线性，因此等于
$\langle z_x, y\rangle_q$，其中 $z_x$ 唯一；记 $u(x) := z_x$，
由唯一性可知其线性，且 $\langle u(x), y\rangle_q =
\varphi'(x,y) = \varphi'(y,x) = \langle u(y), x\rangle_q$：故 $u$
在欧氏空间 $(E,
\langle\cdot,\cdot\rangle_q)$ 中是对称的。谱定理
（\cref{thm:b2:quadratic:spectral}）给出关于 $q$ 标准正交的
特征基 $(\varepsilon_i)$，满足 $u(\varepsilon_i) =
\mu_i\varepsilon_i$：在此基下 $q(x) = \sum x_i^2$ 且 $q'(x) =
\langle u(x), x\rangle_q = \sum\mu_ix_i^2$.

\textbf{21.} 设 $P$ 为该基的矩阵，由合同
可得 $P^{\mathsf T}AP = I$ 且 $P^{\mathsf T}BP =
\operatorname{diag}(\mu_i)$。于是
\[
\det(B - \mu A) = \det(P^{-\mathsf T})
\det\bigl(\operatorname{diag}(\mu_i) - \mu I\bigr)
\det(P^{-1})
= (\det P)^{-2}\prod_i(\mu_i - \mu) :
\]
$\mu_i$ 即矩阵束 $\det(B - \mu A)$ 的根。

\textbf{22.} $\det(B - \mu A) =
\det\begin{pmatrix} -2\mu & 1-\mu\\ 1-\mu & -\mu\end{pmatrix}
= 2\mu^2 - (1-\mu)^2 = \mu^2 + 2\mu - 1$：根为 $\mu_\pm = -1
\pm \sqrt2$。解 $(B - \mu_\pm A)v = 0$，得 $v_\pm = (1 -
\mu_\pm,\ 2\mu_\pm)$（第二行的恒等式 $(1-\mu)^2 =
2\mu^2$ 在根处成立，可验证之）。$A$-范数
很整齐：$q_A(v_\pm) = 2(1 + \mu_\pm^2)$，并验证
$\varphi_A(v_+, v_-) = 0$，利用 $\mu_+ + \mu_- = -2$，
$\mu_+\mu_- = -1$。基 $\Bigl(\frac{v_+}{\sqrt{2(1 +
\mu_+^2)}}, \frac{v_-}{\sqrt{2(1+\mu_-^2)}}\Bigr)$ 构成
关于 $A$ 的标准正交基，且使 $B$ 对角化，对角元为
$\mu_\pm$.

\textbf{23.} 若存在可逆的 $P$ 使两个二次型同时对角化，
则第 21 题的计算将给出 $\det(B - \mu A) = (\det
P)^{-2}\prod(d_{2i} - \mu d_{1i})$，是一个实多项式，
可分解为实一次因式之积。但此处
\[
\det(B - \mu A) = \det\begin{pmatrix} -\mu & 1\\ 1 &
\mu\end{pmatrix} = -\mu^2 - 1 ,
\]
次数为 $2$ 且无实根，矛盾。（
洛伦兹符号差的 $A$ 允许 $B$ 的“旋转”，却无
实轴。）

\textbf{24.} $A^{-1}B = A^{-1/2}\bigl(A^{-1/2}BA^{-1/2}\bigr)
A^{1/2}$，其中 $A^{1/2} = \sqrt A$ 为正定
（见\cref{exo:b2:quadratic:6}）：$A^{-1}B$ 相似于
\emph{对称}矩阵 $A^{-1/2}BA^{-1/2}$，因而可对角化，
且特征值为实数。又 $\det(B - \mu A) = \det A\cdot
\det(A^{-1}B - \mu I)$：第 21 题中矩阵束的根 $\mu_i$
恰为 $A^{-1}B$ 的特征值。

\textbf{25.} (i) 各部分都依赖于谱定理：
经由平方根（楚列斯基分解、极分解）、特征值
界（阿达马不等式）、主轴（二次曲线），以及
适配于 $q$ 的版本（同时化简）。（ii）格拉姆
实现、阿达马不等式与极分解在
半定情形中仍然成立；楚列斯基分解的唯一性、
主元公式以及同时化简都需要定性，
（第 12 题与第 23 题恰好说明它们如何失效）。（iii）
耦合小振动：动能（正定）
与势能是两个二次型；
第 20--22 题中的基即系统的\emph{简正模}，
而 $\mu_i$ 即各角频率的平方。
\end{solution}
